\tocsection{\textbf{9/1/26 Lecture 2: }Topology}

\begin{theorem}[\textbf{Theorem 2.1: }Dynkin's $\pi$-$\lambda$ Theorem]
    If $\cal{C}$ is a $\pi$-system and $\cal{D}$ is a $\lambda$-system such that $\cal{C} \subseteq \cal{D}$, then $\sigma(\cal{C}) \subseteq \cal{D}$. 
\end{theorem}
\begin{proof}{Proof:}
    First, note that as set $S$ is a $\sigma$-algebra if, and only if, it is both a $\pi$ and $\lambda$ system. Clearly, we have the $\Rightarrow$ direction, and for the $\Leftarrow$ direction, note that if $S$ is a $\pi$ and $\lambda$ system, it contains $\Omega, \emptyset$, complements, and countable unions (De Morgan's + $S$ is a $\pi$-system $\implies$ closed under finite unions + increasing limits $\implies$ closed under countable unions).\\
    
    Now, let $\lambda(\cal{C})$ be the smallest $\lambda$-system containing $\cal{C}$ (it is clear that the intersection of $\lambda$-systems is a $\lambda$-system so this definition makes sense). We just need to show that $\lambda(\cal{C})$ is a $\sigma$-algebra, which is to show that it is closed under intersections (making it a $\pi$-system)\\
    
    Let $A, B \in \lambda(\cal{C})$. If both $A, B \in \cal{C}$, then clearly $A \cap B \in \lambda{\cal(C)}$ since $A \cap B \in \cal{C}$. Now, if $B \in \cal{C}$, then consider the set $\cal{A}_{B} = \{A \subseteq \Omega \col A \cap B \in \lambda(\cal{C})\}$. Note that $\cal{A}_{B}$ contains $\cal{C}$. It also contains $\Omega$. Now, if $F \subseteq G$ with $F, G \in \cal{A}_{B}$, since $F \cap B, G \cap B \in \lambda(\cal{C})$ we see that $F \cap B \subseteq G \cap B$, so $(G \cap B) \backslash (F \cap B) = (G \backslash F) \cap B \in \lambda(\cal{C})$, so $G \backslash F \in \cal{A}_{B}$.\\

    Finally, let $A_i \uparrow A$ with $A_i \in \cal{A}_{B}$. Then, as $A_i \cap B \in \lambda(\cal{C})$ and $A_i \cap B \uparrow A \cap B$, we see that $A \cap B \in \lambda(\cal{C})$, so $A \in \cal{A}_{B}$. Thus, $\cal{A}_{B}$ is a $\lambda$-system containing $\cal{C}$, so $\lambda(\cal{C}) \subseteq \cal{A}_{B}$. Thus, we see that if $A \in \lambda{\cal{C}}$ and $B \in \cal{C}$, then $A \cap B \in \lambda(\cal{C})$. \\
    
    Finally, if $A \in \lambda(\cal{C})$, consider the set $\cal{A}_{A}' = \{B \subseteq \Omega \col A \cap B \in \lambda(\cal{C})\}$. By the discussion above, we see that $\cal{C}$ is contained in $\cal{A}_{A}'$, and by the same argument as above, $\cal{A}_{A}'$ is a $\lambda$-system. Thus, $\lambda(\cal{A}) \subseteq \cal{A}_{A}'$, which implies that $A \cap B \in \lambda(\cal{A})$.\\

    Then, as $\sigma(\cal{C}) \subseteq \lambda(\cal{C}) \subseteq \cal{D}$, we have our desired inclusion.

    $\qed$
\end{proof}
This theorem is useless for now, but will soon prove itself useful! Now, some topology.

\begin{definition}[\textbf{Definition 2.2: }Base for a Topology]
    A \emph{base} for a topological space $S$ is a collection of open sets $\cal{C}$ such that every open set in $S$ is a union of sets in $\cal{C}$. Equivalently, a collection $\cal{C}$ is a base if for all open $\cal{U}$ with $x \in \cal{U}$, there exists a $C \in \cal{C}$ such that $x \in C \subseteq \cal{U}$. 
\end{definition}

\begin{definition}[\textbf{Definition 2.3: }Second-Countable Spaces]
    A topological space $S$ is \emph{second-countable} if it has a countable base.
\end{definition}
It can be shown fairly easily that, for metric spaces, second-countability is equivalent to separability, which means that the metric space has a countable dense subset.\\

\begin{definition}[\textbf{Definition 2.4:} Cartesian Product]
    Suppose for some index set $T$, we have sets $X_{t}$, $t \in T$. If $\sqcup_{t \in T} X_{t}$ is the disjoint union of these sets, then the \emph{Cartesian product}, $\bigtimes_{t \in T} X_{t}$ is the set of maps $f$ from $T$ to $\sqcup_{t \in T} X_{t}$ such that $f(t) \in X_{t}$ for all $t \in T$. 
\end{definition}

For example, we see that $\R^2 = \R \times \R$ is the collection of functions $f : \{1, 2\} \to \R \sqcup \R$ such that $f(1) \in \R$ and $f(2) \in \R$. 

\begin{definition}[\textbf{Definition 2.5:} Product Topology / Tychonoff Topology]
    Given topological spaces $S_{t}$, $t \in T$ for some index set $T$, the \emph{product topology} (or \emph{Tychonoff topology}) on $\bigtimes_{t \in T} S_{t}$ is such that open sets in $\bigtimes_{t \in T} S_{t}$ can be written as the union of sets of the form $\bigtimes_{t \in T} \cal{U}_{t}$ where $\cal{U}_{t}$ are open in $S_{t}$ and $\cal{U}_{t} = S_{t}$ for all but finitely many $t \in T$.\\

    In other words, the product topology are generated by sets of the form $\bigtimes_{t \in T} \cal{U}_{t}$ with the properties mentioned.
\end{definition}

One thing a reader may be curious is why we require the finiteness of $\cal{U}_{t} \ne S_{t}$? The topology where we don't require that, called the \emph{box topology} is still a topology, but it isn't as nice. For instance, the product of compact topological spaces is compact under the product topology, but not necessarily under the box topology.\\

Because of this, if $\Omega_{t}$ are topological spaces, then we always endow $\bigtimes_{t} \Omega_{t}$ with the product topology unless otherwise mentioned. 

\begin{lemma}[\textbf{Lemma 2.6}]
    If $S_1, S_2, \ldots$ are second-countable topological spaces, then so is $\bigtimes_{t = 1}^{\infty} S_{t}$. 
\end{lemma}
\begin{proof}{Proof:}
    Let $\cal{C}_{n}$ be a countable base for each $n$. Now, let $\cal{C}$ consist of sets of the form $\bigtimes_{t = 1}^{\infty} C_n$ where $C_n \in \cal{C}_n$ and $C_n = S_n$ for all but finitely many $n$. We see that $\cal{C}$ is countable.\\

    Now, recall that the product topology on $\bigtimes_{t = 1}^{\infty} S_{t}$ has base $\cal{U} = \bigtimes_{t = 1}^{\infty} \cal{U}_t$ where $\cal{U}_{t}$ is open and $\cal{U}_{t} = S_{t}$ for all but finitely $\cal{U}_{t}$. Now, for $(x_1, x_2, \ldots) \in \cal{U}$, since we can find $C_i \in \cal{C}_i$ such that $x_i \in C_i \subseteq \cal{U}_i$, for the finitely many $i$ with $\cal{U}_i \ne S_{i}$, we see that if $C = \bigtimes_{t = 1}^{\infty} C_{t}$, we have $C \in \cal{C}$ and $x \in C \subseteq \cal{U}$. Thus, $\cal{C}$ is also a base for $\bigtimes_{t = 1}^{\infty} S_{t}$, so it has a countable base.

    $\qed$
\end{proof}

\begin{definition}[\textbf{Definition 2.7:} Product $\sigma$-Field]
    Let $\Omega_{t}$, $t \in T$ be measurable spaces, each with a $\sigma$-field $\cal{A}_{t}$. Then, for the Cartesian product, $\bigtimes_{t \in T} \Omega_{t}$, the \emph{product $\sigma$-field} on $\bigtimes_{t \in T} \Omega_{t}$, denoted $\bigotimes_{t \in T} \cal{A}_{t}$ is the $\sigma$-field generated by the \emph{cylinder sets} $A_{t} \times \bigtimes_{s \ne t} \Sigma_{s}$ for $A_{t} \in \cal{A}_{t}$. 
\end{definition}
\pagebreak

\begin{lemma}[\textbf{Lemma 2.7}]
    If $S_1, S_2, \ldots$ are second-countable topological spaces, then
    \[ \cal{B}(S_1 \times S_2 \times \cdots) = \cal{B}(S_1) \otimes \cal{B}(S_2) \otimes \cdots \]
    where $\cal{B}(X)$ gives the Borel $\sigma$-algebra on a topological space $X$. 
\end{lemma}
We will not prove this and just move on. Now, recall that every point map $f : S \to T$ induces a set map $f^{-1} : \P(T) \to \P(S)$ by
\[ f^{-1}(B) = \{s \in S \col f(s) \in B\} \]
Note that $f^{-1}$ commutes with complements, countable unions, and countable intersections; that is
\[ f^{-1}(A^{c}) = (f^{-1}(A))^{c}, \quad f^{-1}\left(\bigcup_{i = 1}^{\infty} A_{i}\right) = \bigcup_{i = 1}^{\infty} f^{-1}(A_i), \quand f^{-1}\left(\bigcap_{i = 1}^{\infty} A_i\right) = \bigcap_{i = 1}^{\infty} f^{-1}(A_i) \]
so that motivates the following lemma.

\begin{lemma}[\textbf{Lemma 2.8}]
    Let $f : S \to T$ where $(S, \cal{S})$ and $(T, \cal{T})$ are measurable spaces. Then,
    \begin{enumerate}
        \item $\cal{S}' = f^{-1}(\cal{T})$ is a $\sigma$-field on $S$.
        \item $\cal{T}' = \{B \subseteq T \col f^{-1}(B) \in \cal{S}\}$ is a $\sigma$-field on $T$. 
    \end{enumerate}
\end{lemma}

